
\noindent
Zur Analyse des Spannungsverhaltens der realen Stromquelle wird die am Lastwiderstand
anliegende Spannung normiert dargestellt. Die relative Spannung wird definiert als:
\begin{equation}
u = \frac{U}{U_0}
\end{equation}

\noindent
Die Spannung am Lastwiderstand ergibt sich aus dem Zusammenhang zwischen Strom
und Leitwert:
\begin{equation}
U = \frac{I}{G_L}
\end{equation}

\noindent
Mit dem Strom durch den Lastleitwert
\begin{equation}
I = I_0 \frac{G_L}{G_L + G_I}
\end{equation}
ergibt sich für die Spannung:
\begin{equation}
U = \frac{I_0}{G_L + G_I}
\end{equation}

\noindent
Als Referenzspannung wird die Leerlaufspannung der Stromquelle $U_0$ verwendet.
Diese ergibt sich für $G_L \rightarrow 0$ zu:
\begin{equation}
U_0 = \frac{I_0}{G_I}
\end{equation}

\noindent
Durch Einsetzen der Ausdrücke für $U$ und $U_0$ in die Definition der relativen
Spannung ergibt sich:
\begin{equation}
u = \frac{G_I}{G_L + G_I}
\end{equation}

\noindent
Zur normierten Darstellung wird erneut das Verhältnis zwischen Lastleitwert und
Innenleitwert eingeführt:
\begin{equation}
r = \frac{G_L}{G_I}
\end{equation}

\noindent
Damit ergibt sich für die relative Spannung die folgende Abhängigkeit vom Leitwert-
verhältnis:
\begin{equation}
u(r) = \frac{1}{1 + r}
\end{equation}
\begin{figure}[H]
    \centering
    \includegraphics[width=1\textwidth]{2_2_4_u_r_stromquelle.png}
    \caption{Strom im Verhältnis der Widerstände (Stromquelle)}
    \label{fig:Strom im Verhaeltnis der Widerstaende (Stromquelle)}
\end{figure}

\noindent
Die relative Spannung nimmt mit zunehmendem Lastleitwert monoton ab. Für sehr
kleine Lastleitwerte ($r \ll 1$) nähert sich die relative Spannung dem Wert eins, wäh-
rend sie für große Lastleitwerte ($r \gg 1$) gegen null strebt.
